\chapter{Improved gauge actions}

There have been a number of methods proposed to improve the gauge
action beyond $O(a^2)$.  These include the Symanzik approach,
mean-field improvement, and renormalization group inspired perfect
actions.  Let me begin this discussion by reviewing the Symanzik
program.

\section{L\"uscher-Weisz Action}

The leading order term in the expansion of all Wilson loops is
$\mathcal{O}^{(4)} = \sum_{\mu\nu} F_{\mu\nu}F_{\mu\nu}$ and corrections
begin at $O(a^2)$ as there are no dimension 5 operators.  Thus, any
lattice action written as a linear combination of Wilson loops will
have the correct continuum limit with corrections at $O(a^2)$.  There
are three dimension 6 operators, which in continuum notation are:
\begin{equation}
\begin{aligned}
\mathcal{O}_1^{(6)} &= \sum_{\mu,\nu} \Tr \bigg(D_\mu F_{\mu\nu}
    D_\mu F_{\mu\nu}\bigg),\\
\mathcal{O}_2^{(6)} &= \sum_{\mu,\nu,\rho} \Tr \bigg(D_\mu F_{\nu\rho}
    D_\mu F_{\nu\rho}\bigg),\\
\mathcal{O}_3^{(6)} &= \sum_{\mu,\nu,\rho} \Tr \bigg(D_\mu F_{\mu\rho}
    D_\nu F_{\nu\rho}\bigg).
\end{aligned}
\label{eq:gaugedimsix}
\end{equation}
Also there are only 3 six-link loops that one can draw on the lattice.
These, as shown in Fig.~\ref{fig:7}, are the planar
$\mathcal{L}_1^{(6)}$, twisted $\mathcal{L}_2^{(6)}$ and the L shaped
$\mathcal{L}_3^{(6)}$ respectively.  Thus, classical improvement of
the lattice action, i.e. removing the $O(a^2)$ term, can be achieved
by taking a linear combination of the plaquette and these three
six-link loops.  Each of these loops has the expansion
\begin{equation}
\mathcal{L} = r^{(4)}\mathcal{O}^{(4)}
    + r_1^{(6)}\mathcal{O}_1^{(6)}
    + r_2^{(6)}\mathcal{O}_2^{(6)}
    + r_3^{(6)}\mathcal{O}_3^{(6)} + \ldots,
\label{eq:gaugeimpa}
\end{equation}

\begin{figure}[htbp]
\centering
\includegraphics[width=0.8\textwidth]{images/fig7.png}
\caption{The three types of six link loops.}
\label{fig:7}
\end{figure}

and L\"uscher and Weisz have shown an elegent way of calculating the
expansion coefficients $r_\alpha^{(d)}$ [73].  Their results are
summarized in Table~\ref{tab:4}.

\begin{table}[htbp]
\centering
\begin{tabular}{lccccc}
\toprule
Loop & $r^{(4)}$ & $r_1^{(6)}$ & $r_2^{(6)}$ & $r_3^{(6)}$\\
\midrule
$\mathcal{L}^{(4)}$  & $-\frac{1}{4}$ & $\frac{1}{24}$ & 0 & 0\\
$\mathcal{L}_1^{(6)}$ & $-2$ & $\frac{5}{6}$ & 0 & 0\\
$\mathcal{L}_2^{(6)}$ & $-2$ & $-\frac{1}{6}$ & $\frac{1}{6}$ & $\frac{1}{6}$\\
$\mathcal{L}_3^{(6)}$ & $-4$ & $\frac{1}{6}$ & 0 & $\frac{1}{2}$\\
\bottomrule
\end{tabular}
\caption{The coefficients of the continuum operators of dimension 4 and
6 in the classical expansion of Wilson loops with 4 and 6 links.}
\label{tab:4}
\end{table}

The lattice gauge action can be written as
\begin{equation}
\mathcal{S}_g = \frac{6}{g^2}\bigg\{c^{(4)}(g^2)\mathcal{L}^{(4)}
    + \sum_{i=1,3} c_i^{(6)}(g^2)\mathcal{L}_i^{(6)}\bigg\}
\label{eq:gaugeimpb}
\end{equation}
in terms of the plaquette and these three 6-link loops.  Using the
results in Table~\ref{tab:4} one gets the normalization condition
\begin{equation}
c^{(4)}(g^2) + 8c_1^{(6)}(g^2) + 8c_2^{(6)}(g^2)
    + 16c_3^{(6)}(g^2) = 1
\label{eq:gaugeimpc}
\end{equation}
such that in the continuum limit the action reduces to
$F_{\mu\nu}F_{\mu\nu}/4$.  Also, from Table~\ref{tab:4} one can see
that tree level improvement can be obtained by the choice
\begin{equation}
c^{(4)} + 20c_1^{(6)} = 0;\quad c_2^{(6)} = 0;\quad c_3^{(6)} = 0.
\label{eq:gaugeclassicalimp}
\end{equation}
This removes $\mathcal{O}_1^{(6)}$ by cancelling the contribution from
the two planar loops.  Of the three dimension six operators,
$\mathcal{O}_1^{(6)}$ alone breaks Lorentz (rotational) invariance.
Thus, rotational invariance of the heavy $q\bar q$ potential can be
used to non-perturbatively tune $c_1^{(6)}$ and to compare the various
improved actions as discussed later.

To include quantum effects, Curci, Menotti, and Paffuti did a 1-loop
calculation.  They imposed the criteria that all Wilson loops be
improved at the leading logarithm level [74,75].  This corresponds to
improving the string tension.  (To kill the $O(a^2)$ terms at all
length scales instead of only in the long distance heavy quark
potential requires including the finite pieces in the 1-loop
calculations.  The fully improved action, as we shall see in the final
L\"uscher-Weisz form, will lead to $O(g^2)$ corrections to $c^{(4)}$
and $c_1^{(6)}$ and the need for $\mathcal{O}_2^{(6)}$.)  The result of
their calculation required the improved action to have the same form as
given in Eq.~\ref{eq:gaugeclassicalimp}.  This agreement with the
classical result is fortuitous.  Symanzik's on-shell improvement
program, in general, prescribes that one has to derive three relations
to fix the $c_i^{(6)}$ such that on-shell quantities do not have
$O(a^2)$ corrections.  L\"uscher and Weisz, in fact, found that
improving all large loops produced only two constraints at tree-level
[76]
\begin{equation}
\begin{aligned}
c^{(4)} + 20c_1^{(6)} - 4c_2^{(6)} + 4c_3^{(6)} &= 0\\
c_2^{(6)} &= 0.
\end{aligned}
\label{eq:gaugeimpd}
\end{equation}
Thus, they were left with a 1-parameter family of improved actions.
The reason for this freedom is the presence of a redundant operator
$\mathcal{O}_3^{(6)}$ generated by the field redefinition
$A_\mu \to A_\mu + a^2g^2d(g)D_\nu F_{\nu\mu}$ [73,76].  Since one has
the liberty to choose the coefficient of an interaction that is part of
a redundant operator without affecting on-shell improvement, the
convenient choice $c_3^{(6)} = 0$ simplifies the action and again leads
us back to Eq.~\ref{eq:gaugeclassicalimp}.  To summarize,
Eq.~\ref{eq:gaugeimpd} are the most general set of constraints that can
be derived based on improving large Wilson loops.  The criteria
employed by Curci et al., improving the large distance part of $q\bar q$
potential, is a subset.  Since this involves a spectral quantity one
can use the redundant operator to get the simpler form
Eq.~\ref{eq:gaugeclassicalimp}.

The relation $c_3^{(6)} = 0$ holds both at tree and 1-loop level as it
comes from exploiting the redundant operator.  Combining this with the
first constraint in Eq.~\ref{eq:gaugeimpd} and the normalization
condition Eq.~\ref{eq:gaugeimpc} leads to a single condition, which at
tree-level is
\begin{equation}
c_1^{(6)} - c_2^{(6)} = -\frac{1}{12},
\end{equation}
and is all that can be gotten from Wilson loops.  To get a second
condition L\"uscher and Weisz had to set up other spectral quantities
to be improved.  Failing to find a suitable condition in the real
theory, they looked at a situation in which the gluons are massive and
appear as asymptotic states.  Using the mass and scattering amplitude
of these states they derived the 1-loop improved action [76], which
goes under the name L\"uscher-Weisz action and has leading correction
at $O(g^4a^4)$,
\begin{equation}
\begin{aligned}
c_0^{(4)}(g^2) &= \frac{5}{3} + 0.2370g^2\\
c_1^{(6)}(g^2) &= -\frac{1}{12} - 0.02521g^2\\
c_2^{(6)}(g^2) &= -0.00441g^2\\
c_3^{(6)}(g^2) &= 0.
\end{aligned}
\end{equation}
In this form for the $c_i$, the normalization condition,
Eq.~\ref{eq:gaugeimpc}, is already incorporated.

The bottom line on the L\"uscher-Weisz action is that since the
coefficient $c_2^{(6)}$ is small, the action, to a very good
approximation, can be improved by keeping just $\mathcal{L}^{(4)}$ and
$\mathcal{L}_1^{(6)}$ as in the classical case.  This 1-loop result
can be mean-field improved without much effort as discussed next.

\section{Tadpole/Mean-field Improvement}

The representation of gauge fields by the exponentiated form
\begin{equation}
U_\mu(x) = e^{iagA_\mu(x)} = 1 + iagA_\mu(x)
    - \frac{a^2g^2}{2}A_\mu^2(x) + \ldots
\end{equation}
gives rise to local quark-gluon vertices with one, two, $\ldots$
gluons.  Of these vertices, all but the lowest are lattice artifacts.
Contracting the two gluons in $a^2g^2A_\mu^2(x)/2$ gives rise to
tadpole diagrams.  Parisi, and Lepage and Mackenzie [77,78] recognized
that these artifacts, though na\"ively suppressed by powers of $a$, are
in fact only suppressed by powers of $g^2$ due to the ultraviolet
divergences generated by the tadpole loops.  They reasoned that these
tadpole contributions were responsible for (i) the poor match between
short distance quantities (like the link in Landau gauge, small Wilson
loops, $\kappa_c$, $\ldots$) and their perturbative estimates, and (ii)
large coefficients in the perturbative lattice expansions.  They,
therefore, proposed a mean-field improvement (also called tadpole
improvement(TI)) to get rid of these artifacts.

The mean-field improvement starts by assuming that the lattice fields
can be split into a UV and IR parts, and that the UV part should be
integrated out [77,78], i.e.
\begin{equation}
e^{iagA_\mu(x)} = e^{iag(A_\mu^{IR}(x)+A_\mu^{UV}(x))}
    = u_0 e^{iagA_\mu^{IR}(x)} \equiv u_0 \tilde U_\mu(x).
\end{equation}
The rescaling of links by an overall constant factor $u_0 < 1$ keeps
the theory gauge invariant.  Under this rescaling (each $U \to u_0
\tilde U$) the Wilson gauge and fermion action becomes
\begin{equation}
\mathcal{S}^W \to \sum_x \bigg[\beta u_0^4\mathcal{W}_{11}
    + \bar{\tilde\psi}\tilde\psi
    - \kappa u_0 \bar{\tilde\psi}\tilde D\tilde\psi\bigg].
\end{equation}
Thus, a simplified recipe for mean-field improvement is to use the
better expansion parameters $\tilde g^2 = g^2/u_0^4$, $\tilde\kappa =
\kappa u_0$, $\tilde U = U/u_0$, and $\tilde\psi =
\psi^{cont}/\sqrt{2\kappa u_0}$.  Note that in terms of the tilde
variables the theory looks the same as the original theory except that
all links are scaled by $u_0$.  Lepage and Mackenzie [78], and
subsequent work, have shown that mean-field improvement of 1-loop
estimates gives, in many cases, results that are very close to
non-perturbative estimates, and the expansion coefficients are much
smaller.  Tadpole improvement is, therefore, a useful thing to do,
especially since it requires negligible extra effort once the 1-loop
calculation has been done.  The only limitation is that it does not
work equally well in all cases and there is no real understanding of
why it does not work well in some cases.  To illustration tadpole
improvement I will apply it to the L\"uscher-Weisz action.

There are two common choices of the tadpole factor $u_0$: (i) the
fourth root of the plaquette, and (ii) the expectation value of the
link in Landau gauge.  So far most calculation have used the first
choice, whereas Lepage and collaborators now advocate the link in
Landau gauge as a better choice [79].  To implement tadpole improvement
requires both the 1-loop expressions for these variables and their
non-perturbative values.  For the Wilson action
\begin{equation}
\begin{aligned}
u_0 &= \langle\mathcal{W}_{11}\rangle^{1/4} = 1 - \frac{1}{12}g^2\\
u_0 &= \langle U\rangle_{LG} = 1 - 0.0772g^2
\end{aligned}
\end{equation}
and for the L\"uscher-Weisz action, Eq.~\ref{eq:gaugeclassicalimp}, it
is [80]
\begin{equation}
u_0 = \langle\mathcal{W}_{11}\rangle^{1/4}
    = 1 - \left(\frac{1}{2}\right)\left(\frac{4}{3}\right)
        \left(\frac{0.366262}{4}\right)g^2
    = 1 - 0.061044g^2.
\end{equation}

Denoting the 1-loop tadpole factor as $u_0 = 1 - Xg^2$, the
tadpole-improved L\"uscher-Weisz (TILW) action is [81]
\begin{equation}
\mathcal{S}_{TILW} = \beta_{eff}\bigg(\mathcal{L}^{(4)}
    - \frac{1}{20u_0^2}\big(1 + (0.1602-2X)g^2\big)\mathcal{L}_1^{(6)}
    - \frac{1}{u_0^2}0.00264g^2\mathcal{L}_2^{(6)}\bigg)
\label{eq:LWTIaction}
\end{equation}
The non-perturbative value of $u_0$ has to be determined
self-consistently as the relative co-efficients in the action
themselves depend on $u_0$.  Note that $\mathcal{L}_2^{(6)}$, which
comes in only as a quantum effect, has, after TI, a coefficient is
reduced to a negligible value.

An important point to note about the TILW action is the strength of
$\mathcal{L}_1^{(6)}$.  Since $u_0 = 0.85-0.90$ for lattice parameters
currently in use, the strength of $\mathcal{L}_1^{(6)}$ is larger than
the classical value $1/20$, i.e. tadpole improvement suggests an
over-correction.  Later on we shall see that this over-correction is
also characteristic of improved actions based on renormalization group
methods.

A simplified description of TI applied to the above example is that the
coefficient of $\mathcal{L}_1^{(6)}$ was multiplied by
$u_0^2(pert)/u_0^2(non-pert)$, the relative factor of two extra links
between the plaquette and $\mathcal{L}_1^{(6)}$.  The factor
$u_0^2(pert)$, which differs from $u_0^2(non-pert)$ at $O(g^4)$, is
absorbed in the 1-loop expansion.  This, as is obvious, makes the
coefficient of $g^2$ small.  The correction is, therefore, dominated by
$u_0(non-pert)$ which can be measured with very high precision.  There,
however, is the residual uncertainty due to the freedom in the choice of
$u_0$.  To summarize, the underlying philosophy of TI is that the large
tadpole contributions are present in the perturbative expansion of all
quantities, and can be absorbed by factoring $u_0(pert)$ raised to the
appropriate power depending on the number of links.  This makes the
expansion coefficients small.  The modified perturbative expansion is
then in terms of an improved $\tilde g^2 = g^2/u_0^4$, and the
correction is dominated by overall factors of $u_0(non-pert)$, which
should be chosen to be the expectation value of the link in Landau
gauge.

One of the first tests to convincingly demonstrate improvement using
the TILW action was restoration of rotational invariance in the $q\bar
q$ potential.  The results from [82] are reproduced in Fig.~\ref{fig:8}.
Since then tests of improvement using the scaling of the finite
temperature transition, $T_c$, and topological susceptibility have been
carried out.  These results, as will be discussed later, all make a
strong case for improvement.

\begin{figure}[htbp]
\centering
\includegraphics[width=0.8\textwidth]{images/fig8.png}
\caption{Static-quark potential computed on $6^4$ lattices with
$a \approx 0.4$ fm using the Wilson action and the TILW action.  The
dotted line is the standard infrared parameterization for the continuum
potential, $V(r) = Kr - \pi/12r + c$, adjusted to fit the on-axis
values of the potential.  Figure reproduced from [82].}
\label{fig:8}
\end{figure}

\section{The Renormalization Group and the Renormalized Trajectory}

This is a brief summary of Wilson's formulation of the renormalization
group and the concept of the renormalized trajectory.  For details see
[83--90].  I will start with a discussion in momentum space.  Consider
the effective theory formulated at some scale $Q$ in terms of the
generalized action
\begin{equation}
\mathcal{S}(Q) = \sum_\alpha K_\alpha^{(0)}(Q)S_\alpha
\end{equation}
where $S_\alpha$ are all possible operators (interactions) and
$K_\alpha^{(0)}$ are the corresponding coupling constants.  Now
consider an exact renormalization group transformation (RGT)
$\mathcal{T}_b$ that integrates all momenta between $Q$ and $Q/b \equiv
Q/2$.  Here $b = 2$ is the scale factor of the transformation.  The
renormalized theory at scale $Q/2$ is
\begin{equation}
\mathcal{S}(Q/2) = \sum_\alpha K_\alpha^{(1)}(Q/2)S_\alpha
\end{equation}
where $K_\alpha^{(1)}$ are analytical functions of $K_\alpha^{(0)}$,
\begin{equation}
K_\alpha^{(1)} = \mathcal{T}_b\big[K_\alpha^{(0)}\big].
\end{equation}
The sequence of theories defined by repeated use of $\mathcal{T}_b$
defines a flow trajectory in the space of couplings.  The physical
theory is preserved under each $\mathcal{T}_b$, and thus along the whole
trajectory.  Correlation lengths, measured in lattice units, change
trivially $\xi/a \to \xi/ab$.  Fixed points of this transformation are
defined by
\begin{equation}
K_\alpha^{(*)} = \mathcal{T}_b\big[K_\alpha^{(*)}\big],
\end{equation}
i.e. points at which the theory reproduces itself under
$\mathcal{T}_b$.  Since all correlation lengths (momenta) scale by $b$
under $\mathcal{T}_b$, therefore at the fixed points
\begin{equation}
\xi\big|_{\mathrm{fixed\ point}} = 0 \quad \mathrm{or} \quad \infty.
\end{equation}
Fixed points with $\xi = 0$ are trivial fixed points.  Examples of
these are the $T = 0$ and $T = \infty$ limit of most statistical
mechanics models.  We are interested in non-trival fixed points with
$\xi = \infty$.  At these points the correlation length $\xi/a =
1/Ma$, measured in units of lattice spacing, is infinite.  This allows
us to take the continuum limit of lattice field theories --- set
$a = 0$ holding $M$ fixed.

Now consider a point $\{K_\beta\}$ close to the fixed point
$\{K_\beta^*\}$ and a RGT with an infinitesimal $b$.  Under the
assumption that there are no singularities in the space of coupling
constants, the new couplings $K_\alpha'(K_\beta)$ generated by
$\mathcal{T}_b$ can be written as a Taylor expansion about the fixed
point
\begin{equation}
K_\alpha'(K_\beta) = K_\alpha^* +
    \frac{\partial K_\alpha'}{\partial K_\beta}\bigg|_{K^*}\Delta K_\beta
    + \ldots.
\label{eq:ltma}
\end{equation}
In that case one can define a ``linear region'' in the neighborhood of
$\{K_\beta^*\}$ in which
\begin{equation}
\Delta K_\alpha' =
    \frac{\partial K_\alpha'}{\partial K_\beta}\bigg|_{K^*}\Delta K_\beta
\label{eq:ltmb}
\end{equation}
where $\Delta K_\beta = K_\beta - K_\beta^*$.  The linearized
transformation matrix (also called the stability matrix)
\begin{equation}
\mathcal{T}_{\alpha\beta} \equiv
    \frac{\partial K_\alpha'}{\partial K_\beta}\bigg|_{K^*}
\label{eq:ltmc}
\end{equation}
controls the RG flows locally.  In field theory discussions, the
stability matrix $\mathcal{T}_{\alpha\beta}$ is called the
$\beta$-function.  Note that, because of the truncation of
Eq.~\ref{eq:ltma} all statements based on the analysis of
$\mathcal{T}_{\alpha\beta}$ have the implicit assumption that they are
approximately true in the linear region close to the fixed point and
for small $b$, and become exact only as $\{K_\beta\} \to
\{K_\beta^*\}$.

Wilson formulated the real space renormalization group as an alternate
to the momentum space renormalization group discussed above.  This
approach is characterized by a constrained integral over the position
dependent original degrees of freedom $s$ to produce an effective
theory,
\begin{equation}
e^{K_\alpha'S_\alpha(s')} = \int \mathcal{D}s\,\mathcal{T}(s',s)
    e^{K_\alpha S_\alpha(s)}.
\end{equation}
Here $\mathcal{T}(s',s)$ is the constraint that specifies how the $s$
in some local neighborhood (called the block cell) are to be averaged
and replaced by an effective $s'$.  A necessary condition on
$\mathcal{T}(s',s)$, that preserves the partition function, is
\begin{equation}
\int \mathcal{D}s'\,\mathcal{T}(s',s) = 1.
\label{eq:blockcondition}
\end{equation}
One can see that if the block cell in $D$-dimensions is $b^D$, and the
constraint $\mathcal{T}(s',s)$ is equivalent to integrating momenta in
the range $Q/2 - Q$, then the two approaches are equivalent.  In
general, an exact integration of the upper half momenta is not feasible
and $\mathcal{T}(s',s)$ is constructed as some sensible ``average'' of
the original spins in the block cell.  Other than the condition
Eq.~\ref{eq:blockcondition}, the definition of a sensible ``average''
is that it preserves the physical properties of the theory.  Because of
this dual description, momentum versus real space formulations, the RGT
is also called a blocking transformation --- a usage more common in
statistical mechanics and numerical implementations.

A brief summary of the properties of the matrix
$\mathcal{T}_{\alpha\beta}$, evaluated using either the momentum or
real space formulations, and its associated eigenvalues $\Lambda_i$ and
eigenvectors $\Phi_i$ obtained by solving $\mathcal{T}\Phi_i =
\Lambda_i\Phi_i$ is as follows.

\begin{enumerate}
\item $\mathcal{T}_{\alpha\beta}$ is, in general, not a symmetric
  matrix.  Its eigenvalues are, however, real.

\item The eigenvalues depend on the scale factor $b$.  Since the theory
  obtained after two successive transformations, $b \otimes b$, should
  be the same as the one obtained after a single scale change by $b^2$,
  the general form of the eigenvalues is $\Lambda_i = b^{\lambda_i}$
  where the $\lambda_i$ are independent of $\mathcal{T}_b$ and are
  called the critical exponents.

\item The eigenvectors $\Phi_i$ are by definition the linear scaling
  fields.  Consider a small deviation of size $u$ from the fixed point
  along a given eigenvector $\Phi$.  Then under a RGT, the renormalized
  value of the coupling, $u'$, is given by
  \begin{equation}
  u'\Phi = \mathcal{T}_b u\Phi = b^{\lambda_\Phi}u\Phi.
  \label{eq:ltmd}
  \end{equation}
  Note the crucial use of the assumption of linearity in deducing this
  simple scaling of $u$.  The eigenvectors depend on the precise form of
  $\mathcal{T}_b$.

\item Eigenvalues $\Lambda_i > 1$ lead to flows away from the fixed
  point as deviations along these $\Phi_i$ grow under $\mathcal{T}_b$.
  These eigenvalues and the corresponding eigenvectors are called
  relevant.

\item Operators corresponding to eigenvalues $\Lambda_i < 1$ die out
  geometrically with the number of blocking steps.  These operators do
  not contribute to the long distance properties of the system and are
  consequently called irrelevant.  The associated exponents $\lambda_i$
  control corrections to scaling.

\item Eigenvalues that are exactly unity are called marginal.
  Deviations along marginal directions do not change under RGT.  To
  ascertain that an eigenvalue is truly marginal, one has to go beyond
  the linear approximation.  Typically marginal operators develop
  logarithmic corrections at higher order.  A quick examination of the
  $\beta$-function of asymptotically free field theories should convince
  you that $g$ is a marginal coupling at tree-level and becomes relevant
  due to quantum corrections.

\item There is an additional class of eigenvectors, called redundant
  operators, that are not physical.  Redundant operators are
  manifestation of the invariance of the action under a redefinition of
  the fields, i.e., they define relations between the different
  operators such as the equations of motion.  The corresponding
  eigenvalues can be $<1,\ =1,$ or $>1$.  These eigenvalues depend on
  the choice of the RGT, i.e. different RGT with the same or different
  scale factor $b$ will give different $\lambda_i$.  Since different
  non-linear RGT correspond to different rescaling of the fields,
  redundant operators are specific to the transformation.  In fact, one
  way to isolate them is to repeat the calculation of
  $\mathcal{T}_{\alpha\beta}$ with a different RGT; the exponents
  $\lambda_i$ that change with the RGT correspond to redundant
  operators.  Even though they carry no physics information it is
  important to isolate these operators, especially in numerical
  calculations.  For example, when improving the action one wants to use
  all possible redundant operators to simplify the action and thus
  minimize the computational cost.

\item Under certain conditions $\mathcal{T}_{\alpha\beta}$ is block
  diagonal, i.e. the couplings break up into sub-sets that are closed
  under the RGT.  For example, the spin flip symmetry $s \to -s$ of the
  Ising Hamiltonian at zero magnetic field causes the interactions to
  break up into odd and even sectors.  One can, therefore, calculate
  $\mathcal{T}_{\alpha\beta}$ for each sector separately as the
  eigenvalue spectrum is independent.
\end{enumerate}

In this generalized space of couplings $\{K_\alpha\}$, there exist
points, called critical points $\{K_\alpha^c\}$, with $\xi = \infty$.
Starting from a critical point, a RGT produces another since $\xi' =
\xi/b = \infty$.  Thus the set of critical points define a hypersurface
in the infinite dimensional space $\{K_\alpha\}$.  The RG flows on this
surface can (a) meander randomly, (b) go to some limit cycle or a
strange attractor, or (c) converge to a fixed point $\{K_\alpha^*\}$.
We shall only be concerned with possibility (c).  At the fixed points
the renormalized couplings are exactly equal to the original couplings,
$K_\alpha' = K_\alpha$, and the theory reproduces itself at all length
scales.  This is in distinction to critical points where only the long
distance behavior is reproduced.  The location of the fixed point
depends on the RGT and in general a given RGT will have more than one
fixed point with $\xi = \infty$.  Each such fixed point has a basin of
attraction i.e. the set of $\{K_\alpha^c\}$ that converge to it under
the RGT.  This hypersurface is thus orthogonal to all relevant
eigenvectors of $\mathcal{T}_b$.  The basin of attraction defines the
universality class since the long distance behavior of all theories is
the same and governed by the fixed point theory.

The flow along a relevant scaling field, whose end-point is the fixed
point is called the renormalized trajectory (RT).  Just as the fixed
point is an attractor of all critical points in its universality class,
the RT is the attractor of all flows terminating in the critical points
that lie in the basin of attraction of the fixed point.  This is
illustrated in Fig.~\ref{fig:9}.  For simplicity I have assumed that
there is only one relevant coupling, as is the case for QCD.  Along the
RT there are no scaling violations as, by construction, all the
irrelevant scaling fields are zero.  Thus, simulations done using an
action along the exact RT, will reproduce the continuum physics without
discretization errors.  The RT is thus the holy grail of attempts to
improve the lattice formulation, i.e. the perfect action.

\begin{figure}[htbp]
\centering
\includegraphics[width=0.8\textwidth]{images/fig9.png}
\caption{A schematic of the critical hypersurface ($\xi = \infty$), the
fixed point, the renormalized trajectory (RT) flowing towards the
infrared under blocking transformations, and flows attracted by them.
Also shown is the simple Wilson action and flows starting from it.}
\label{fig:9}
\end{figure}

Simulations on finite lattices are done in a region of coupling space
where all $\xi \ll L$, and not at the fixed point or at the critical
points.  Current typical values lie in the range $\xi/a = 1/\sqrt{\sigma}a
= 5$--$15$.  Even for quenched simulations this range will not change in
the near future.  Points in this region lie on flows that terminate at
critical points that may lie quite far from the fixed point.
Corrections to scaling along these flows at $\xi \approx 10$ may,
therefore, be large.  One can reduce these corrections by adjusting the
action to lie on another flow that starts closer to the fixed point and
at $\xi \approx 10$ is closer to the RT.  Thus, to improve the action a
non-perturbative approach is to estimate the RT.  Since flows are
attracted by the RT, its location can be estimated by starting with even
the simple plaquette action and following the flows generated under
some blocking transformation.  Evaluating the blocked action along the
flow is therefore the basis of all renormalization group methods for
improving the action.  Methods used for determining the blocked action
are reviewed in [86].

A standard assumption underlying the validity of the renormalization
group is that the fixed point action is local (the strength of the
couplings fall off exponentially with the size of the Wilson loops and
the separation between points that are coupled by the Dirac operator),
nevertheless, a local action still involves an ``infinite'' number of
couplings.  Keeping all these when evaluating the action generated
under blocking, and in simulations, is impractical.  So one is forced to
to truncate the action to a finite (small) number of interactions.
Also, the RT at $\xi \approx 10$ may be sufficiently different in form
from that near the fixed point.  Consequently, the questions one has to
face when designing improved actions are (i) what is the minimal set of
terms to keep, (ii) are all the redundant interactions removed from this
set, and (iii) how to compare such improved actions evaluated using
different RGT.  Also note that the projection of the RT on to a truncated
subspace may not be the action with the smallest scaling violations that
one can design in that space.  I shall explore attempts to answer these
questions next.

I will summarize current activity in designing improved gauge actions
by reviewing those that are being used to calculate phenomenologically
relevant quantities.  Unfortunately, there is no standard notation used
by different collaborations.  To facilitate comparison I will therefore
try to write them in a common form as well as maintain contact with the
original notation.  Wilson loops will be denoted by $W^i_{\mu\nu}$
where the index $i$ will be used for both the type of loop and its
representation.  Unless otherwise stated, it will be implicit that the
real part of the trace of Wilson loops is to be taken.  The
$W^i_{\mu\nu}$ are normalized such that the maximum value of the trace
is unity.  As before $\beta$ is the overall coupling and the $c_i$ give
the relative strengths of different loops.

\section{Iwasaki Action}

Iwasaki studied block transformations in perturbation theory.  His
criterion was to obtain a gauge action that after a few blockings comes
close to the RT.  His Ans\"atz is [91]
\begin{equation}
\mathcal{S}_{Iwasaki} = \beta\sum_{\mu,\nu}\sum_x
    \bigg(W^{1\times1}_{\mu\nu} - 0.0907W^{1\times2}_{\mu\nu}\bigg)
\label{eq:iwasakiaction}
\end{equation}
where $W^{1\times2}$ is the rectangular loop $\mathcal{L}_1^{(6)}$
shown in Fig.~\ref{fig:7}.  (Iwasaki used the notation
$(1-8c_0)W^{1\times1}_{\mu\nu} + c_0W^{1\times2}_{\mu\nu}$ with
$c_0 = -0.331$ which is the same as Eq.~\ref{eq:iwasakiaction} with a
rescaled $\beta$.)  This action shows better rotational invariance and
scaling of $T_c/\sqrt{\sigma}$.  It is being pursued by the JLQCD
collaboration [92].

\section{$b=\sqrt{3}$ blocking transformation}

The $\sqrt{3}$ blocking transformation has certain very nice properties
for 4-dimensional gauge theories in addition to being a smaller scaling
factor than the conventional $b = 2$ [89].  The block cell consists of a
site and its eight neighbors.  As a result the transformation has much
higher rotational symmetry for both fermions and gauge fields [89].
Gupta and Patel have studied the flow under blocking using Monte Carlo
Renormalization Group methods.  They investigated the effect of (i)
truncation and (ii) variation in the $c_i$ as a function of $\beta$
[87].  For phenomenological studies they approximate the RT by a linear
trajectory in a four parameter space
\begin{equation}
\begin{aligned}
\mathcal{S}_{\sqrt{3}} = \beta\,\mathrm{Re}\Big\{&
    c_F^{1\times1}\sum\Tr W^{(1\times1)}
    + c_F^{1\times2}\sum\Tr W^{(1\times2)}\\
    &+ c_6^{1\times1}\sum\Big[\frac{3}{2}(\Tr W^{(1\times1)})^2
        - \frac{1}{2}\Tr W^{(1\times1)}\Big]\\
    &+ c_8^{1\times1}\sum\Big[\frac{9}{8}|\Tr W^{(1\times1)}|^2
        - \frac{1}{8}\Big]\Big\},
\end{aligned}
\label{eq:rootthreedef}
\end{equation}
with
\begin{equation}
c_6^{1\times1}/c_F^{1\times1} = -0.12;\quad
c_8^{1\times1}/c_F^{1\times1} = -0.12;\quad
c_F^{1\times2}/c_F^{1\times1} = -0.04.
\label{eq:rootthreerat}
\end{equation}
This particular choice of couplings (truncation of the action) was made
for the following reasons.  The coupling for the plaquette in the
adjoint representation, $c_8^{1\times1}$, was included to avoid the
phase structure in the fundamental-adjoint plane (see Section 15); that
in the 6-dimensional representation, $c_6^{1\times1}$, since it differs
from $c_8^{1\times1}$ only at $O(a^4)$ and higher, and the $1\times2$
rectangle in the fundamental representation, $c_F^{1\times2}$, as it is
the simplest loop that kills the $O(a^2)$ terms a la Symanzik.

Classical expansion, see Eq. 7.3, of the loops gives the normalization
condition
\begin{equation}
\bigg[c_F^{1\times1} + \frac{5}{2}c_6^{1\times1}
    + \frac{9}{4}c_8^{1\times1}\bigg] + 8c_F^{1\times2} = 1,
\end{equation}
where I have lumped together the first three terms coming from
$W^{(1\times1)}$ as the effective contribution of the plaquette.  Using
the ratios given in Eq.~\ref{eq:rootthreerat}, one gets a very high
weight $c^{(4)} = 1/0.11$, which as shown in [51] is correlated with a
suppression of dislocations versus physical instantons.  Similarly, the
coefficient of the leading discretization correction
$\mathcal{O}_1^{(6)}$, using Table~\ref{tab:4}, is
\begin{equation}
\frac{1}{24}\bigg[c_F^{1\times1} + \frac{5}{2}c_6^{1\times1}
    + \frac{9}{4}c_8^{1\times1}\bigg]
    + \frac{5}{6}c_F^{1\times2}.
\end{equation}
For the ratios given in Eq.~\ref{eq:rootthreerat} this is not zero, but
over-corrected by the proposed strength of $W^{(1\times2)}$.  At
$1/a \approx 2$ GeV, this over-correction is larger than that for the
tadpole improved L\"uscher-Weisz Ans\"atz, similar to
$\mathcal{S}_{Iwasaki}$, and smaller than that in
$\mathcal{S}_{QCDTARO}$.

Early spectrum studies and measurement of the glueball spectrum [93] to
test improvement were inconclusive.  In hindsight the failure to
quantify improvement is due to poor statistics and not looking at
quantities like rotational invariance at strong coupling, where the
lattice artifacts show up clearly for the simple Wilson action.
Recently this action has also been used to study lattice topology and
the results show that short distance artifacts are suppressed [51].

\section{$b=2$ blocking transformation}

The QCDTARO collaboration has also investigated the flow of the action
under blocking [94,95].  They use a $b = 2$ blocking transformation and
evaluate the couplings on blocked lattices using Schwinger-Dyson
equations.  Their overall conclusion is that the blocked action, to a
very good approximation, can be represented by the two couplings
$W^{1\times1}$ and $W^{1\times2}$.  Their estimate of the RT cannot be
approximated by a linear relation between the coefficients.  An
approximate Ans\"atz in the range of couplings simulated is
\begin{equation}
\mathcal{S}_{QCDTARO} = \beta\sum_i\sum_{\mu,\nu}\sum_x
    \bigg(W^{1\times1}_{\mu\nu} - 0.11W^{1\times2}_{\mu\nu}\bigg)
\end{equation}
Note that the coefficient of $W^{1\times2}$ is larger than the tadpole
improved Luscher-Weisz, Iwasaki, or the $b = \sqrt{3}$ actions.

\section{Classically Perfect Action}

Hasenfratz and Niedermayer showed that one can derive the classically
perfect action for asymptotically free-field theories using a
saddle-point integration of the RGT about $g = 0$ [96,97].  They use the
$b = 2$ blocking transformation proposed by Swendsen [98].  The
practical issue here is whether the saddle-point integration gives an
action that can be approximated by a few local terms.  The application
of this method to SU(3) and the tuning of the truncated version is
still under investigation [99].  The current version of the truncated
classical fixed-point action (CFP) is [100]
\begin{equation}
\mathcal{S}_{CFP} = \beta\sum_{i,x}\{c_1^i(1 - \mathrm{Re}\,\Tr W^i)
    + \ldots + c_4^i(1 - \mathrm{Re}\,\Tr W^i)^4\}.
\label{eq:Perfectaction}
\end{equation}
where the sum over $i$ involves two kinds of loops, the plaquette and
the 6-link loop that twists around a 3-cube ($\mathcal{L}_2^{(6)}$ in
Fig.~\ref{fig:7}).  Higher powers of the loops arise naturally in the
saddle-point integration.  Since the strengths decrease rapidly, the
authors choose to truncate at the fourth power.  The choice of the
particular 6-link loop is motivated by the desire to keep the action
local in time and to help improve rotational invariance.  The values of
coefficients $c_1$--$c_4$ for the two types of loops are given in
Table~\ref{tab:5}.

\begin{table}[htbp]
\centering
\begin{tabular}{lccccc}
\toprule
Loop & $c_1$ & $c_2$ & $c_3$ & $c_4$\\
\midrule
$1\times1$        & 3.248   & $-1.580$ & 0.1257  & 0.0576\\
Twisted 6 link    & $-0.2810$ & 0.0051 & 0.0049  & $-0.0096$\\
\bottomrule
\end{tabular}
\caption{The current coefficients of the truncated ``perfect'' action
with 4 and 6 link loops.}
\label{tab:5}
\end{table}

At this point full control over stability with respect to truncation
has not been achieved.  The second issue is whether this Ans\"atz,
obtained at $a = 0$, continues to be a useful parameterization at
$1/a \sim 2$ GeV where simulations are done, i.e. is the renormalized
trajectory approximately linear.  To test this requires a
non-perturbative tuning about $1/a \sim 2$ GeV using the same $b = 2$
blocking transformation used in the saddle-point integration.  The
calculations necessary to do this tuning have not yet been initiated.
As a result this Ans\"atz is considered as work in progress.
Nevertheless, the tests for improvement in the scaling of $T_c$, string
tension, topological susceptibility, and in the restoration of
rotational symmetry are already very encouraging [100,52].

\section{Summary}

The question --- which is the most improved gauge action --- is still
far from settled.  As should be obvious, the drawback of the
renormalization group improved actions is that the expansion is in the
size of the loops or the dimensionality of the representation and not
in the dimensionality of the irrelevant operators.  Since each loop
contributes to operators of all dimensions it is not clear how to
parameterize the residual error ($O(a^2) + O(a^4)$, or $O(a^4)$, or
higher).  Also, implementing the equations of motions (removing the
freedom associated with redundant directions) has not been easy.  So
there is not yet a consensus on the obvious minimal set of operators
that should be kept.  Different operators have been added based on (i)
their importance in perturbative expansions, (ii) size of the
coefficients in blocked actions (keep in mind that the freedom of the
choice of the blocking transformation and consequently the associated
RT effects these coefficients), and (iii) specific incorporation of
non-perturbative effects; for example the inclusion of a negative
adjoint coupling helps avoid the phase structure in the
fundamental-adjoint plane (see Section 15), and adding loops mimicking
the $F^{\mu\nu}\tilde F_{\mu\nu}$ operator help give scale invariant
instantons [100,101].  One hopes that a more precise method for tuning
the gauge action will emerge in the future.  Fortunately, in spite of
this ambiguity, the situation is very encouraging.  The improvement
seen in the restoration of rotational invariance (measured using $q\bar
q$ potential as shown in Fig.~\ref{fig:8}), glueball spectrum, and
finite temperature studies is remarkable.  To a good approximation,
most of these actions are very well represented by adding just the
$W^{(1\times2)}$ term with a coefficient larger than the value suggested
by tadpole improvement, and keeping a few higher representations of the
plaquette.  Such a truncated space keeps the computational problem
tractable.

The bottom line is that even though there is not a commonly accepted
improved gauge action, the improvements seen with any of the above
actions are significant.  At present the more pressing challenge, in
terms of the size of the errors and the computational cost, is improving
the Dirac discretization.  This is discussed next.
